\section*{Резюме}
\addcontentsline{toc}{section}{Резюме}
\label{sec:rezyume}

Курсовият проект изследва ролята на прокурора в досъдебното производство по
Наказателно-процесуалния кодекс на Република България. Работата започва с
конституционната и законовата основа на прокурорските правомощия (чл.~127 КРБ, чл.~46
НПК). След това проследява правомощията на прокурора при предварителната проверка и
образуването на производството, при ръководството и надзора върху разследването, при
мерките за процесуална принуда и при приключването на досъдебната фаза. Специално
внимание е отделено на границата между правомощията, които прокурорът упражнява сам, и
тези, за които е необходимо разрешение на съда. Отделна част разглежда особеностите при
разследването на киберпрестъпления: запазването и изземването на компютърни данни,
достъпа до данните за трафика, изземването на криптоактиви и трансграничното събиране на
електронни доказателства, включително по Регламент (ЕС) 2023/1543.

Изводът е, че прокурорът е централният орган на досъдебното производство и отговаря за
неговия резултат, но най-тежките намеси в правата на гражданите остават под
предварителен съдебен контрол. При киберпрестъпленията значението на прокурора нараства,
защото бързината на неговите решения определя дали електронните доказателства ще бъдат
запазени.

\medskip
\noindent\textbf{Ключови думи:} прокурор; досъдебно производство; ръководство и надзор;
съдебен контрол; киберпрестъпления; електронни доказателства.
